% GENERATED FILE. Edit canonical lessons, then run npm run generate:books.
% canonical-source-hash: fnv1a64:aa895720af1bd024
% canonical-lessons: LA-C392-caveo, LA-C392-cautio, LA-C392-necessitas, LA-C392-opportunus, LA-C392-praecaveo

\chapter{Taking Care}
\label{ch:la-taking-care}
\hlchaptermodality{392}

\begin{chapteropening}
\textbf{By the end of this chapter:} \emph{I can say that I am on my guard and take precautions, and talk about a guarantee, a necessity and a good moment.}

Complete the last lesson of chapter 392: I can say that I am on my guard and take precautions, and talk about a guarantee, a necessity and a good moment.
\end{chapteropening}

\section[caveō, cavēre]{caveō, cavēre — to beware, to be on one's guard}

\label{lesson:LA-C392-caveo}



\begin{quote}
Before the new one: say the Latin for differently, then the Latin for or if.
\end{quote}

\subsection*{You'll want to know: caveō, cavēre}
\textbf{caveō, cavēre} — \textquotedblleft{}to beware, to be on one's guard\textquotedblright{}.

\begin{quote}
In mūrō scrīptum erat: CAVĒ CANEM.
\end{quote}

\begin{quote}
\emph{On the wall was written: BEWARE OF THE DOG.}
\end{quote}

\emph{caveō, cavēre, cāvī, cautum}. \emph{Cavēre} + accusative, to beware of; \emph{cavē nē} + subjunctive, take care not to. English \emph{caution} and \emph{caveat}, let him beware.

\begin{cousinweb}
\emph{caveō} is inherited from Proto-Italic \emph{*kawēō} (EtymDB), from a root meaning \textquotedblleft{}to watch, to heed\textquotedblright{}. English \emph{caution} and \emph{caveat} come from it.
\end{cousinweb}

\begin{grammarlens}[title={What you've built}]
One more word for telling, explaining and describing at B1.
\end{grammarlens}

\subsection*{Your turn}
\begin{itemize}
\raggedright
  \item \emph{Say it:} \emph{caveō, cavēre}
  \item \emph{Say it:} \emph{caveō, cavēre}, once more
  \item \emph{Say it:} a sentence of your own with \emph{caveō, cavēre}
  \item \emph{Recall:} say the Latin for differently, then the Latin for or if, then say \emph{caveō, cavēre} again
\end{itemize}

\begin{tcolorbox}[breakable,colback=teal!4,colframe=teal!35!black,title={Before you move on}]
What does caveō, cavēre mean? (\textquotedblleft{}To beware, to be on one's guard\textquotedblright{}.) Say it once more.
\end{tcolorbox}
\section[cautiō, cautiōnis]{cautiō, cautiōnis — caution; a guarantee, a bond}

\label{lesson:LA-C392-cautio}



\begin{quote}
Before the new one: say the Latin for or if, then the Latin for to beware.
\end{quote}

\subsection*{You'll want to know: cautiō, cautiōnis}
\textbf{cautiō, cautiōnis} — \textquotedblleft{}caution; a guarantee, a bond\textquotedblright{}.

\begin{quote}
Sine cautiōne pecūniam nōn crēdit.
\end{quote}

\begin{quote}
\emph{He does not lend money without a guarantee.}
\end{quote}

\emph{cautiō, cautiōnis} f., from \emph{caveō}: carefulness, and in law a written security. English \emph{caution}.

\begin{cousinweb}
\emph{cautiō} is the noun built on \emph{caveō} \textquotedblleft{}to beware, to take care\textquotedblright{} (EtymDB; Lewis \& Short). English \emph{caution} comes from it.
\end{cousinweb}

\begin{grammarlens}[title={What you've built}]
One more word for telling, explaining and describing at B1.
\end{grammarlens}

\subsection*{Your turn}
\begin{itemize}
\raggedright
  \item \emph{Say it:} \emph{cautiō, cautiōnis}
  \item \emph{Say it:} \emph{cautiō, cautiōnis}, once more
  \item \emph{Say it:} a sentence of your own with \emph{cautiō, cautiōnis}
  \item \emph{Recall:} say the Latin for or if, then the Latin for to beware, then say \emph{cautiō, cautiōnis} again
\end{itemize}

\begin{tcolorbox}[breakable,colback=teal!4,colframe=teal!35!black,title={Before you move on}]
What does cautiō, cautiōnis mean? (\textquotedblleft{}Caution; a guarantee, a bond\textquotedblright{}.) Say it once more.
\end{tcolorbox}
\section[necessitās, necessitātis]{necessitās, necessitātis — necessity, need; compulsion}

\label{lesson:LA-C392-necessitas}



\begin{quote}
Before the new one: say the Latin for to beware, then the Latin for caution.
\end{quote}

\subsection*{You'll want to know: necessitās, necessitātis}
\textbf{necessitās, necessitātis} — \textquotedblleft{}necessity, need; compulsion\textquotedblright{}.

\begin{quote}
Necessitās mē coēgit.
\end{quote}

\begin{quote}
\emph{Necessity forced me.}
\end{quote}

\emph{necessitās, necessitātis} f., the noun of \emph{necesse}, needful. \emph{Necessitās nōn habet lēgem}, necessity knows no law. English \emph{necessity}.

\begin{cousinweb}
\emph{necessitās} is built on \emph{necesse} \textquotedblleft{}needful\textquotedblright{} (Lewis \& Short), which EtymDB takes as \emph{ne-} \textquotedblleft{}not\textquotedblright{} with \emph{cessus}, from \emph{cēdō} \textquotedblleft{}to give way\textquotedblright{}: what does not give way. English \emph{necessity} comes from it.
\end{cousinweb}

\begin{grammarlens}[title={What you've built}]
One more word for telling, explaining and describing at B1.
\end{grammarlens}

\subsection*{Your turn}
\begin{itemize}
\raggedright
  \item \emph{Say it:} \emph{necessitās, necessitātis}
  \item \emph{Say it:} \emph{necessitās, necessitātis}, once more
  \item \emph{Say it:} a sentence of your own with \emph{necessitās, necessitātis}
  \item \emph{Recall:} say the Latin for to beware, then the Latin for caution, then say \emph{necessitās, necessitātis} again
\end{itemize}

\begin{tcolorbox}[breakable,colback=teal!4,colframe=teal!35!black,title={Before you move on}]
What does necessitās, necessitātis mean? (\textquotedblleft{}Necessity, need; compulsion\textquotedblright{}.) Say it once more.
\end{tcolorbox}
\section[opportūnus, opportūna, opportūnum]{opportūnus, opportūna, opportūnum — convenient, timely, opportune}

\label{lesson:LA-C392-opportunus}



\begin{quote}
Before the new one: say the Latin for caution, then the Latin for necessity.
\end{quote}

\subsection*{You'll want to know: opportūnus, opportūna, opportūnum}
\textbf{opportūnus, opportūna, opportūnum} — \textquotedblleft{}convenient, timely, opportune\textquotedblright{}.

\begin{quote}
Tempus opportūnum exspectāmus.
\end{quote}

\begin{quote}
\emph{We are waiting for a convenient time.}
\end{quote}

\emph{opportūnus, opportūna, opportūnum}: said first of a wind blowing towards the harbour, then of anything well timed. English \emph{opportune}.

\begin{cousinweb}
\emph{opportūnus} is built from \emph{ob} \textquotedblleft{}towards\textquotedblright{} and \emph{portus} \textquotedblleft{}a harbour\textquotedblright{} (Lewis \& Short): of a wind that drives you into port. English \emph{opportune} and \emph{opportunity} come from it.
\end{cousinweb}

\begin{grammarlens}[title={What you've built}]
One more word for telling, explaining and describing at B1.
\end{grammarlens}

\subsection*{Your turn}
\begin{itemize}
\raggedright
  \item \emph{Say it:} \emph{opportūnus, opportūna, opportūnum}
  \item \emph{Say it:} \emph{opportūnus, opportūna, opportūnum}, once more
  \item \emph{Say it:} a sentence of your own with \emph{opportūnus, opportūna, opportūnum}
  \item \emph{Recall:} say the Latin for caution, then the Latin for necessity, then say \emph{opportūnus, opportūna, opportūnum} again
\end{itemize}

\begin{tcolorbox}[breakable,colback=teal!4,colframe=teal!35!black,title={Before you move on}]
What does opportūnus, opportūna, opportūnum mean? (\textquotedblleft{}Convenient, timely, opportune\textquotedblright{}.) Say it once more.
\end{tcolorbox}
\section[praecaveō, praecavēre]{praecaveō, praecavēre — to take precautions, to guard against beforehand}

\label{lesson:LA-C392-praecaveo}



\begin{quote}
Before the new one: say the Latin for necessity, then the Latin for convenient.
\end{quote}

\subsection*{You'll want to know: praecaveō, praecavēre}
\textbf{praecaveō, praecavēre} — \textquotedblleft{}to take precautions, to guard against beforehand\textquotedblright{}.

\begin{quote}
Ab īnsidiīs praecāvit.
\end{quote}

\begin{quote}
\emph{He took precautions against the ambush.}
\end{quote}

\emph{praecaveō, praecavēre, praecāvī, praecautum}: \emph{prae-} beforehand + \emph{caveō}. English \emph{precaution}.

\begin{cousinweb}
\emph{praecaveō} is built from \emph{prae-} \textquotedblleft{}before\textquotedblright{} and \emph{caveō} \textquotedblleft{}to beware\textquotedblright{} (EtymDB; Lewis \& Short). English \emph{precaution} comes from its family.
\end{cousinweb}

\begin{grammarlens}[title={What you've built}]
One more word for telling, explaining and describing at B1.
\end{grammarlens}

\subsection*{Your turn}
\begin{itemize}
\raggedright
  \item \emph{Say it:} \emph{praecaveō, praecavēre}
  \item \emph{Say it:} \emph{praecaveō, praecavēre}, once more
  \item \emph{Say it:} a sentence of your own with \emph{praecaveō, praecavēre}
  \item \emph{Recall:} say the Latin for necessity, then the Latin for convenient, then say \emph{praecaveō, praecavēre} again
\end{itemize}

\begin{tcolorbox}[breakable,colback=teal!4,colframe=teal!35!black,title={Before you move on}]
What does praecaveō, praecavēre mean? (\textquotedblleft{}To take precautions, to guard against beforehand\textquotedblright{}.) Say it once more.
\end{tcolorbox}
